\chapter{Conception décisionnelle}
\label{ch:decisionnel}

\section{Introduction}
Ce chapitre détaille la conception décisionnelle de la solution~: les sources de données exploitées, l'architecture du processus ETL, sa mise en œuvre avec Talend, et la conception de l'entrepôt de données qui constitue le socle de l'ensemble de la chaîne de restitution.

\section{Description des sources de données}
La principale source de données exploitée est l'\textbf{Enquête Ménages-Déplacements (EMD)} de Dakar, complétée par des données de \textbf{comptage de trafic} issues de capteurs et de relevés sur les principaux sites de comptage du réseau. L'EMD est structurée en plusieurs blocs de variables, identifiées par des codes (I2, I4, I9, M21, M26, M55, M66 à M91, etc.), couvrant~:
\begin{itemize}
\item les caractéristiques géographiques du domicile (quartier, strate, district de recensement)~;
\item les caractéristiques du ménage (taille, type de logement, équipement)~;
\item les déplacements individuels (motif, mode, durée, coût)~;
\item l'accessibilité aux transports en commun et aux services (école, daara)~;
\item les problèmes et manquements perçus dans le quartier (manque de transport, inondations, enclavement, insécurité).
\end{itemize}

Ces données sources ont fait l'objet d'une priorisation en trois phases afin de structurer le développement de l'ETL de façon incrémentale~:
\begin{enumerate}
\item \textbf{Phase 1 -- MVP}~: colonnes géographiques (I2, I4, I9), accessibilité de base (M66 à M70), problèmes majeurs (M73, M87) et taille du ménage (M21)~;
\item \textbf{Phase 2 -- Extension services}~: variables d'accès aux écoles et daaras (M90, M91)~;
\item \textbf{Phase 3 -- Détails avancés}~: variables de résilience en cas de pluie (M69), modes de transport vers les services, et autres problèmes de quartier (bruit, insécurité, accidents).
\end{enumerate}

\section{Dictionnaire de données et qualité des sources}

\subsection{Dictionnaire de données}
Afin de documenter précisément le contenu des fichiers sources et leur correspondance avec le modèle décisionnel, un \textbf{dictionnaire de données} a été constitué (fichier \texttt{Dictionnaire\_Donnees\_CETUD.xlsx}, dossier \texttt{CETUD DATA}). Il couvre les quatre tables sources exploitées~: \texttt{Menage} (231 colonnes brutes), \texttt{Individu} (210 colonnes brutes), \texttt{Deplacement} (11 colonnes) et \texttt{Trafic} (19 colonnes), ainsi qu'un onglet \texttt{Mapping\_DW} reliant explicitement chaque colonne source retenue à sa dimension ou son indicateur de destination dans l'entrepôt. Le tableau~\ref{tab:mapping_dw} reprend la synthèse de ce mapping.

\begin{table}[H]
\centering
\caption{Mapping entre les tables de faits, leurs dimensions et leurs tables sources EMD}
\label{tab:mapping_dw}
\begin{tabular}{p{2.6cm}p{3.4cm}p{2.4cm}p{5.2cm}}
\toprule
\textbf{Table de faits} & \textbf{Dimensions liées} & \textbf{Source EMD} & \textbf{Mesures principales} \\
\midrule
Fait\_Comptage & Dim\_Site, Dim\_Temps & Trafic & nb\_vehicules, volume horaire, flux entrant/sortant, vitesse moyenne \\
Fait\_Deplacement & Dim\_Motif, Dim\_Mode, Dim\_Geographie & Deplacement & durée trajet, heure départ/arrivée, coût total \\
Fait\_IndividuMenage & Dim\_Individu, Dim\_Menage, Dim\_Geographie & Individu + Menage & nb déplacements/jour, taux équipement voiture, taux motorisation \\
\bottomrule
\end{tabular}
\end{table}

\begin{figure}[H]
\centering
\includegraphics[width=0.85\textwidth]{diction1.png}
\caption{Capture d'écran -- page de présentation du dictionnaire de données CETUD}
\label{fig:dico_presentation}
\end{figure}

\begin{figure}[H]
\centering
\includegraphics[width=0.85\textwidth]{diction2.png}
\caption{Capture d'écran -- onglet Mapping\_DW du dictionnaire de données}
\label{fig:dico_mapping}
\end{figure}

\subsection{Qualité des données sources}
Les données exploitées sont des \textbf{données réelles, fournies par le CETUD} (et non des données simulées), issues de l'Enquête Ménages-Déplacements et des campagnes de comptage de trafic. Une phase de \emph{data discovery} a été menée sous \textbf{Google Colab} (notebook \texttt{EDA\_CETUD\_DataDiscovery.ipynb}, dossier \texttt{CETUD DATA/02\_eda\_notebook}) afin d'évaluer la qualité de ces sources avant leur intégration dans l'ETL. Cette analyse exploratoire a mis en évidence~:
\begin{itemize}
\item un \textbf{volume de colonnes brutes très hétérogène} entre tables (jusqu'à 231 colonnes pour \texttt{Menage}), nécessitant une phase de sélection/priorisation des variables réellement exploitables (cf. priorisation en trois phases ci-dessus)~;
\item une \textbf{proportion significative de valeurs manquantes} sur certaines variables déclaratives (notamment le taux de motorisation des ménages et la mobilité des seniors), illustrée à la figure~\ref{fig:eda_nulls}~;
\item des \textbf{formats hétérogènes} entre les exports Excel bruts (\texttt{.xls}/\texttt{.xlsx}) et les fichiers CSV transformés en entrée de Talend, imposant une normalisation des types avant chargement~;
\item des \textbf{codes de variables non auto-documentés} (codes EMD du type \texttt{P1}, \texttt{M21}, \texttt{M66}, etc.), justifiant la construction du dictionnaire de données comme référentiel central de traduction.
\end{itemize}

\begin{figure}[H]
\centering
\includegraphics[width=0.75\textwidth]{eda_nulls_distribution.png}
\caption{Data discovery sous Google Colab -- distribution des valeurs manquantes par variable}
\label{fig:eda_nulls}
\end{figure}

Ce travail de qualification des sources, réalisé en amont de l'ETL, a directement conditionné la priorisation des variables intégrées (section~\ref{ch:decisionnel} ci-dessus) et explique l'absence de fiabilité suffisante pour certains indicateurs (taux de motorisation, mobilité des seniors), signalée également au chapitre~\ref{ch:resultats}.

\subsection{Choix du modèle décisionnel}
La modélisation multidimensionnelle de l'entrepôt s'appuie sur un \textbf{schéma en constellation}, c'est-à-dire un modèle comportant plusieurs tables de faits partageant un ensemble commun de dimensions. Ce choix, plutôt qu'un schéma en étoile classique limité à une seule table de faits, est justifié par la nature même des données du projet~: les données de comptage de trafic (issues de capteurs, granularité site/heure) et les données d'enquête ménages-déplacements (granularité individu/déplacement) correspondent à des processus métier distincts, avec des cycles de mise à jour différents, mais partagent des axes d'analyse communs (temps, géographie, site). Un schéma en étoile unique aurait nécessité de forcer ces deux processus métier dans une seule table de faits, au prix d'une perte de granularité ou de la création de mesures non applicables selon le contexte~; le schéma en constellation permet au contraire de conserver une table de faits dédiée à chaque processus, tout en mutualisant les dimensions partagées.

\begin{figure}[H]
\centering
\includegraphics[width=0.85\textwidth]{modele_constellation_v2.png}
\caption{Schéma en constellation du modèle décisionnel (4 faits et 9 dimensions)}
\label{fig:constellation}
\end{figure}

\section{Architecture ETL}
La chaîne ETL est organisée en deux niveaux successifs~:
\begin{enumerate}
\item une \textbf{zone de préparation (Staging Area, SA)}~: jobs chargeant les données brutes, préalablement nettoyées sous Python/Colab, dans des tables intermédiaires, sans transformation lourde~;
\item une \textbf{zone de chargement final}~: jobs lisant les tables de la Staging Area, appliquant les transformations, jointures et calculs nécessaires (via le composant \texttt{tMap}), puis chargeant les dimensions et les faits dans l'entrepôt de données (DW).
\end{enumerate}

Cette architecture à deux niveaux a été retenue afin de séparer clairement le chargement brut (rapide, peu de logique) de la transformation (jointures de référentiel, résolution de clés étrangères, calculs de mesures), ce qui facilite le débogage et la traçabilité de chaque étape.

\section{Processus ETL avec Talend}

\subsection{Extraction}
L'extraction est réalisée par les composants \texttt{tFileInputExcel} et \texttt{tFileInputDelimited}, qui lisent les fichiers sources (extraits EMD et comptages de trafic), déjà nettoyés en amont via des scripts Python (traitement des valeurs manquantes, normalisation des types, suppression des doublons).

\subsection{Nettoyage}
Au sein de Talend, un nettoyage complémentaire est appliqué via le composant \texttt{tFilterColumns}, qui sélectionne uniquement les colonnes utiles conformément à la priorisation définie en section~\ref{ch:decisionnel} ci-dessus, réduisant ainsi la complexité des jobs et le volume de données transporté.

\subsection{Transformation}
La transformation constitue le cœur du processus ETL, portée par le composant \texttt{tMap}. Trois jobs principaux assurent d'abord l'alimentation de la zone de préparation~: \texttt{SA\_individu}, \texttt{SA\_menage} et \texttt{SA\_deplacement}, correspondant aux trois grandes entités sources de l'EMD.

\begin{figure}[H]
\centering
\includegraphics[width=0.85\textwidth]{talend/SA_deplacement.png}
\caption{Job Talend de chargement de la Staging Area -- déplacements}
\label{fig:sa_deplacement}
\end{figure}

Le job \texttt{Job\_deplacement} illustre ensuite la logique de transformation appliquée pour construire la table de faits \texttt{Fact\_deplacement}. Le composant \texttt{tMap} central réalise les opérations suivantes~:
\begin{itemize}
\item récupération du flux principal contenant les attributs du déplacement (numéro de ménage, numéro d'individu, lieux et heures de départ/arrivée, durée, motif, fréquence de transport en commun, mode, coût total)~;
\item trois jointures (\emph{lookup}) sur les dimensions \texttt{Dim\_Motif} (par \texttt{motif}), \texttt{Dim\_Individu} (par clé composite \texttt{num\_individu + num\_menage}) et \texttt{Dim\_Mode} (par \texttt{mode}) afin de résoudre les clés étrangères~;
\item construction du flux de sortie combinant les clés étrangères résolues et les mesures (durée du trajet, lieux, coût total, heures de départ/arrivée, fréquence de transport en commun).
\end{itemize}

\begin{figure}[H]
\centering
\includegraphics[width=\textwidth]{talend/tmap_depl.png}
\caption{Composant tMap du job de construction de \texttt{Fact\_deplacement}}
\label{fig:tmap_depl}
\end{figure}

Le job de construction de \texttt{Fact\_Comptage} suit une logique similaire~: une entrée principale portant les mesures de comptage est enrichie par deux jointures de référence via un composant \texttt{tMap}, avant chargement final. Les jobs \texttt{fact\_accessibilite} et \texttt{fact\_indivmen} suivent le même schéma général, en exploitant respectivement les variables d'accessibilité (durée de marche, modes disponibles, exposition aux inondations) et les variables croisées individu/ménage (taille du ménage, équipement, motorisation).

\begin{figure}[H]
\centering
\includegraphics[width=0.85\textwidth]{talend/fact_compt.png}
\caption{Job Talend de construction de \texttt{Fact\_Comptage}}
\label{fig:fact_compt}
\end{figure}

\subsection{Chargement}
Le chargement final dans l'entrepôt de données est réalisé par le composant \texttt{tDBOutput}, qui écrit successivement les dimensions de référence (\texttt{Dim\_Site}, \texttt{Dim\_Geographie}, \texttt{Dim\_Individu}, \texttt{Dim\_Menage}) puis les tables de faits, afin de garantir l'intégrité référentielle au moment du chargement.

\begin{figure}[H]
\centering
\begin{subfigure}[b]{0.48\textwidth}
\includegraphics[width=\textwidth]{talend/config_dim_site.png}
\caption{Configuration de \texttt{Dim\_Site}}
\end{subfigure}
\hfill
\begin{subfigure}[b]{0.48\textwidth}
\includegraphics[width=\textwidth]{talend/dim_geo.png}
\caption{Chargement de \texttt{Dim\_Geographie}}
\end{subfigure}
\caption{Jobs de construction des dimensions de référence}
\end{figure}

\section{Conception du Data Warehouse}

\subsection{Tables de faits}
Le modèle retenu comporte trois tables de faits.

\paragraph{Fait\_Comptage.} Enregistre les mesures issues des dispositifs de comptage automatique de trafic, à la granularité du site et de l'heure (clés étrangères~: \texttt{FK\_id\_Site}, \texttt{FK\_id\_temps}~; mesures~: nombre de véhicules, volume horaire matin/soir, flux entrant, flux sortant, désagrégation directionnelle).

\paragraph{Fait\_Deplacement.} Décrit chaque déplacement individuel déclaré, à la granularité du déplacement (clés étrangères~: \texttt{FK\_id\_motif}, \texttt{FK\_id\_mode}, \texttt{FK\_id\_zone\_residence}~; mesures~: durée du trajet, heures de départ/arrivée, coût total, temps de parcours par motif).

\paragraph{Fait\_IndividuMenage.} Croise les informations individuelles et celles du ménage, à la granularité de l'individu enquêté (clés étrangères~: \texttt{FK\_id\_individu}, \texttt{FK\_id\_menage}, \texttt{FK\_id\_zone\_residence}~; mesures~: nombre de déplacements journaliers, taux d'équipement en voiture, taux de motorisation, mobilité par zone).

\subsection{Dimensions}
Le modèle comporte neuf dimensions, dont le dictionnaire est présenté au tableau~\ref{tab:dim_dict}.

\begin{table}[H]
\centering
\caption{Dictionnaire des dimensions du modèle décisionnel}
\label{tab:dim_dict}
\begin{tabular}{p{2.6cm}p{6.8cm}p{3.6cm}}
\toprule
\textbf{Dimension} & \textbf{Attributs principaux} & \textbf{Granularité} \\
\midrule
Dim\_Site & id\_Site, code\_site, Ouvl\_Soc, catégorie\_unité, DESC\_Site, Sens & Site de comptage \\
Dim\_Temps & id\_temps, année, mois, jour, heure, heure début/fin déplacement & Heure / jour \\
Dim\_Motif & id\_motif, code\_motif, motif\_origine & Motif de déplacement \\
Dim\_Individu & id\_individu, num\_individu, sexe, statut matrimonial, profession, permis & Individu \\
Dim\_Mode & id\_mode, code\_mode, libellé, catégorie\_mode & Mode de transport \\
Dim\_Menage & id\_menage, num\_menage, statut occupation, type logement, équipements & Ménage \\
Dim\_Geographie & id\_zone\_residence, code\_zone, quartier, strate, département & Zone géographique (507 zones) \\
Dim\_Service & id\_service, code\_service, libell\'e\_service, type\_service, op\'erateur & Service de transport \\
Dim\_Difficulte & id\_difficulte, code\_difficulte, libell\'e\_difficulte, cat\'egorie\_difficulte & Type de difficult\'e d'acc\`es \\
\bottomrule
\end{tabular}
\end{table}

\subsection{Hiérarchies}
Plusieurs hiérarchies d'agrégation sont exploitables sur ces dimensions~: \texttt{Dim\_Temps} (heure $\rightarrow$ jour $\rightarrow$ mois $\rightarrow$ année), \texttt{Dim\_Geographie} (zone de résidence $\rightarrow$ quartier $\rightarrow$ strate $\rightarrow$ département), et \texttt{Dim\_Individu} (individu $\rightarrow$ ménage). Ces hiérarchies permettent de restituer les indicateurs à différents niveaux de granularité dans le cockpit Knowage, sans nécessiter de retraitement de l'entrepôt.

\subsection{Mesures et indicateurs}
Au-delà des mesures directement issues des sources, plusieurs indicateurs calculés ont été intégrés au modèle~: le \textbf{score d'accessibilité} (0 à 100), composite, intégrant une pénalité liée à la durée de marche jusqu'à l'arrêt le plus proche (jusqu'à $-40$ points), un bonus de diversité des modes disponibles (jusqu'à $+20$ points), une pénalité de zone inondable ($-25$ points, aggravée de $-20$ points supplémentaires en l'absence de mode disponible en cas de pluie) et une pénalité de tarification abusive en période de pluie ($-15$ points)~; la \textbf{vulnérabilité climatique}, combinant exposition aux inondations et absence de solution de repli~; et l'\textbf{impact sur l'éducation}, mesurant le lien entre accessibilité aux transports et accès aux structures scolaires. Le score d'accessibilité permet de classer chaque zone selon une typologie à quatre niveaux~: \textit{Bien desservie} ($\geq 75$), \textit{Moyenne} ($\geq 50$), \textit{Sous-desservie} ($\geq 25$) et \textit{Critique} ($< 25$).

\subsection{Implémentation physique}
L'entrepôt de données a été déployé dans un schéma dédié (\texttt{DW}) au sein d'une base PostgreSQL. Chaque table de dimension est préfixée \texttt{Dim\_} et chaque table de faits \texttt{Fact\_}, conformément aux conventions de nommage usuelles en modélisation décisionnelle. Les clés primaires des dimensions sont des entiers auto-incrémentés (\texttt{PK\_id\_...}), repris en clé étrangère dans les tables de faits (\texttt{FK\_id\_...}). La structure physique de chaque table a été vérifiée via pgAdmin, par interrogation du catalogue système \texttt{information\_schema.columns}.

\begin{figure}[H]
\centering
\includegraphics[width=0.55\textwidth]{pg/Fact_comptage.png}
\caption{Structure de la table \texttt{Fact\_Comptage} (pgAdmin)}
\label{fig:fact_comptage_pg}
\end{figure}

La dimension géographique (\texttt{Dim\_Geographie}) illustre la richesse du référentiel mis en place~: elle recense 507 zones de résidence, chacune associée à une strate et un quartier de domicile. Une table \texttt{fact\_access} a par ailleurs été ajoutée au modèle initial afin de porter spécifiquement les mesures liées à l'accessibilité (score d'accessibilité, durée de marche, nombre de modes disponibles, indicateur de zone inondable), en lien avec les dimensions \texttt{Dim\_Difficulte}, \texttt{Dim\_Service} et \texttt{Dim\_Transport}.

\begin{figure}[H]
\centering
\includegraphics[width=0.6\textwidth]{pg/dim_geo.png}
\caption{Extrait de la table \texttt{Dim\_Geographie} (507 zones)}
\label{fig:dim_geo_pg}
\end{figure}

L'outil pgAdmin a été utilisé tout au long du projet pour vérifier la structure des tables après chaque chargement Talend, contrôler les volumes chargés, détecter les valeurs manquantes ou aberrantes, et exécuter des requêtes SQL de contrôle qualité. Ce contrôle a notamment permis d'identifier l'absence de données fiables pour le taux de motorisation des ménages et la mobilité des seniors, indicateurs qui n'ont donc pu être que partiellement renseignés.



\section{Conclusion}
Ce chapitre a détaillé la conception décisionnelle de la solution~: les sources de données, l'architecture et la mise en œuvre du processus ETL avec Talend, et la conception complète de l'entrepôt de données PostgreSQL. Ce socle décisionnel alimente directement le cockpit Knowage et l'ensemble des restitutions présentées au chapitre~\ref{ch:resultats}. Le chapitre suivant détaille le développement de la plateforme applicative qui expose cette chaîne décisionnelle aux utilisateurs finaux.

\clearpage
